
\subsection{Sell part of a recovered lot without losing its remaining units}
The whole-lot sale above leaves no recovered stock to track. A company that
sells only part of that stock needs a quantity and cost record for what remains.
The existing lot already records both; its carrying amount alone would not
identify how many units were sold. A separately selected method therefore
takes the sold quantity as an explicit input rather than deriving it from cash
proceeds or inventory value.

This method assumes that the remaining units are homogeneous and divisible,
that their measurement basis has not changed, and that proportional allocation
is the selected treatment. The caller also supplies affirmative findings for
completed performance and the retained equality of tax base and carrying
amount. Those findings record explicit premises. Unknown findings remain
unresolved, and false findings place the event outside this method's scope.
A heterogeneous lot or a different tax treatment requires another method.

Immediately before a sale, let $N>0$ be the remaining units, $C\geq0$ their
carrying amount in USD, and $T=C$ their tax base under the selected treatment.
The event supplies sold units $q$, with $0<q\leq N$, and cash proceeds $P\geq0$
in USD. The fraction applies to the stock immediately before this event:
\[
 f=\frac{q}{N},\qquad
 C_{\mathrm{released}}=fC,\qquad
 T_{\mathrm{released}}=fT.
\]
The remaining lot keeps its identity and acquisition date, with units $N-q$,
carrying amount $C-fC$ and tax base $T-fT$. It becomes disposed exactly when
$q=N$. A positive-quantity lot whose carrying amount is zero remains live;
zero cost does not imply that no goods remain.

The journal debits cash by $P$ and cost of sales by $fC$, and credits revenue
by $P$ and inventory by $fC$. Its debit-positive sum is
$P+fC-P-fC=0$. Cash rises by $P$ and inventory falls by $fC$, so net assets and
profit both rise by $P-fC$. In the indirect operating cash calculation, adding
back the inventory reduction to that profit effect gives
$(P-fC)+fC=P$, the same receipt as the direct cash record. The reduction in
the lot's carrying amount equals the inventory posting; every other lot and
the separate inventory balance remain unchanged. This sale assesses no tax
expense and makes no new measurement of inventory.

After the February settlement in the shared-book example, the recovered lot
has two units with carrying amount and tax base of 20. Cash is 452 and total
inventory is 80. Selling one unit for 15 releases 10 of cost, leaves one unit
with carrying amount and tax base of 10, and gives cash 467 and inventory 70.
The prior lease gain of 20, cost-of-sales credit of 4 and handling expense of 8
give recorded profit of 16 before this sale; its margin of 5 raises profit to
21. Operating cash since the opening is $-48+15=-33$. Assets of 827 less
liabilities of 214 equal equity of 613, also opening equity of 592 plus profit
of 21.

If the remaining unit is then sold for another 15 under the same conditions,
the new fraction is one: one unit is sold from one remaining unit. The second
sale releases the remaining 10, leaves cash 482 and total inventory 60, and
marks the lot disposed. The combined account changes equal those of the
original whole-lot sale for 30. The two event identities and cash receipts
remain separate; equal aggregate effects do not make their histories identical.

A proposal retains the entire before-state, event, allocation and resulting
book. Application recomputes the complete proposal. A changed remaining
quantity, altered allocation or balanced but different journal fails that
comparison. Selling more than the available units, selling a disposed lot,
reusing an event identity or backdating a sale is refused without changing
the successful preceding state.

This optional transition has checked quantity, journal and statement identities
and reproduces the original full-sale state when all units are sold. The
existing whole-lot method remains unchanged. Connecting the optional event to
the mixed-book authoring and annual dispatch interfaces is separate work;
these local controls do not establish that those interfaces already support it.
